\documentclass[12pt]{beamer}
\usetheme{Pittsburgh}
\usecolortheme{seagull}
%\usepackage[utf8]{inputenc}
\usepackage{fontspec}
\usepackage[english]{babel}
\usepackage{amsmath}
\usepackage{amsfonts}
\usepackage{amssymb}
\usepackage{graphicx}
\author{Apresentação e Programa da Disciplina}
\title{Estruturas de Dados}
%\setbeamercovered{transparent}
\setbeamertemplate{navigation symbols}{}
%\logo{\includegraphics[scale=0.015]{Brasao_UFSC.png}\includegraphics[scale=0.2]{brasao_PPGCC.jpg}}
\institute{Departamento de Informática e de Estatística \\ Prof. Jean Everson Martina \\ Prof. Aldo von Wangenheim}
\date{2016.2}
\subject{}
\usebackgroundtemplate{\includegraphics[width=\paperwidth,
height=\paperheight]{../reusable_images/fundo_UFSC.png}}
\begin{document}

{
\usebackgroundtemplate{\includegraphics[width=\paperwidth,
height=\paperheight]{../reusable_images/fundo_capa.png}}
\begin{frame}
\titlepage
\includegraphics[scale=0.3]{../reusable_images/brasao_INE.png}
\end{frame}
}

\begin{frame}{Organização da Disciplina}
\begin{itemize}
\item Objetivos
\item Metodologia
\item Avaliação
\item Material Didático
\item Realização dos Trabalhos
\item Programa da Disciplina
\item Bibliografia
\item Contato com os Docentes
\end{itemize}
\end{frame}

\begin{frame}{Objetivos da Disciplina}
\begin{itemize}
\item Aprender as técnicas básicas de organização de dados em um computador sob a forma de estruturas de dados:
\begin{itemize}
\item na memória;
\item em meios externos.
\end{itemize}
\item Aprender a manipular e implementar estruturas de dados utilizando diversas técnicas de programação atualmente em uso.
\item Aprender a projetar e a avaliar a eficiência de algoritmos complexos.
\end{itemize}
\end{frame}


\begin{frame}{Objetivos da Disciplina}
O aluno deverá:
\begin{enumerate}
\item Adquirir o conhecimento necessário para desenvolver programas de computador eficazes e eficientes para a solução de problemas complexos.
\item Adquirir a experiência prática de desenvolvimento de programas de computador para desenvolver e estimar adequadamente a dificuldade de desenvolvimento de programas complexos.
\end{enumerate}
\begin{block} {Importante!}
O domínio de Estruturas de Dados será indispensável para praticamente tudo o que você for aprender daqui para frente.
\end{block}
\end{frame}

\begin{frame}{Metodologia de Ensino}
\begin{itemize}
\item Disciplina eminentemente prática, com ênfase na aplicação dos conhecimentos;
\item Aulas teóricas alternadas com aulas práticas;
\begin{itemize}
\item Serão vistos aspectos de modelagem, algoritmos, técnicas de implementação e bases teóricas dos modelos.
\end{itemize}
\item Aulas práticas em laboratório:
\begin{itemize}
\item Ministradas pelo estagiário de docência e pelo professor;
\item Sempre iniciarão com uma complementação visando a implementação do que foi visto na aula teórica;
\item Objetivam prover oportunidade para o aluno fixar o que foi visto e aprimorar as suas técnicas de programação.
\end{itemize}
\end{itemize}
\end{frame}

\begin{frame}{Metodologia de Ensino - Implementação}
\begin{itemize}
\item Metodologias de Implementação/Codificação:
\begin{itemize}
\item Serão vistas em paralelo às estruturas de dados;
\item Reveremos a evolução histórica das técnicas atualmente em uso;
\end{itemize}
\item Técnicas de Implementação de Programas:
\begin{itemize}
\item Programação Orientada a Objetos;
\item Meta-Programação;
\item Uso de recursos avançados da Linguagem;
\end{itemize}
\item Linguagens de Programação adequadas às Técnicas:
\begin{itemize}
\item Discussão da adequação das linguagens;
\item Exemplos em outras linguagens.
\end{itemize}
\end{itemize}
\end{frame}

\begin{frame}{Critérios de Avaliação}
\begin{itemize}
\item 7 Notas (Pesos: 2,1,1,1,2,1,2):
\begin{itemize}
\item 2: média dos pequenos trabalhos (entre 12 e 18);
\item 1: projeto de Implementação I;
\item 1: teste parcial teórico ;
\item 1: teste parcial prático de Implementação;
\item 2: projeto de Implementação II;
\item 1: teste final prático de Implementação;
\item 2: teste final teórico ;
\end{itemize}
\end{itemize}
\begin{block}{Avaliação das Implementações:}
Automatizada pelo moodle segundo critérios definidos no enunciado
\end{block}
\end{frame}

\begin{frame}{Recuperação}
\begin{itemize}
\item Conceito ''Menção I" somente será cogitado de acordo com a legislação da UFSC;
\item A recuperação pode não ocorrer porque a disciplina é mais que 50\% prática;
\item A recuperação será realizada através de um trabalho de implementação seguido de entrevista;
\begin{itemize}
\item O trabalho de recuperação será individual;
\item O trabalho de recuperação terá peso igual ao da média do semestre, de acordo com a legislação da UFSC.
\end{itemize}
\end{itemize}
\end{frame}

\begin{frame}{Material Didático}
\begin{itemize}
\item Material da disciplina disponível na Internet (http://moodle.ufsc.br/):
\begin{itemize}
\item Fio condutor das aulas para o aluno;
\item Não usar como única fonte;
\item Cronogramas, listas de trabalhos e datas;
\item Plano de Ensino;
\item Transparências dadas em aula;
\item Arquivos de dados para exercícios e programas-exemplo;
\item Links para sites com manuais e documentação útil para programação.
\end{itemize}
\item Livros.
\end{itemize}
\end{frame}

\begin{frame}{Realização dos Trabalhos}
\begin{itemize}
\item Os trabalhos pequenos são todos individuais;
\item Equipes de 2 pessoas (somente trabalhos grandes):
\begin{itemize}
\item É importante aprender a realizar trabalhos de implementação em equipe;
\end{itemize}
\item Implementação:
\begin{itemize}
\item ANSI C++ - OOP;
\item Plataforma homologada é somente o VPL;
\end{itemize}
\item Ferramentas \& Máquinas:
\begin{itemize}
\item Aproveite as Aulas de Laboratório;
\item Linux e Eclipse para C/C++;
\item Compiladores: gcc / g++
\end{itemize}
\end{itemize}
\end{frame}

\begin{frame}{Entrega dos Trabalhos}
\begin{itemize}
\item Data:
\begin{itemize}
\item Deverão ser entregues na data especificada no Moodle;
\item São aceitos com atraso, mas com 50\% de desconto na nota;
\item Projetos de implementação também;
\end{itemize}
\item Serão aceitos os trabalhos que:
\begin{itemize}
\item Satisfizerem as especificações / enunciados dados;
\item Utilizarem a linguagem e a técnica de programação especificadas;
\item Utilizarem codificação de acordo com os padrões internacionais para a linguagem especificada;
\item Estiverem devidamente documentados.
\item Executarem na plataforma homologada.
\end{itemize}
\end{itemize}
\end{frame}

\begin{frame}{Controle de Freqüência}
\begin{itemize}
\item O controle de freqüência segue as regulamentações da UFSC para cursos presenciais;
\item É obrigatório 75\% de presença para aprovação;
\item O controle de freqüência vai ser feito diretamente no moodle da disciplina;
\item O controle de freqüência poderá ser feito de forma automatizada usando a carteirinha do RU;
\begin{itemize}
\item Se você ainda não tem uma é so ir na secretaria do RU e emitir a sua.
\end{itemize}
\end{itemize}
\end{frame}

\begin{frame}{Programa da Disciplina}
\begin{itemize}
\item Introdução e Aspectos Gerais da Disciplina;
\item Introdução à Programação C++;
\item Listas, Pilhas e Filas utilizando Vetores (arrays);
\item Gerência e Alocação Dinâmica de Memória;
\item Lista Encadeada;
\item Fila e Pilha Baseados em Listas Encadeadas;
\item Lista Duplamente Encadeada;
\item Lista Circular;
\end{itemize}
\end{frame}

\begin{frame}{Programa da Disciplina}
\begin{itemize}
\item Conceitos de Complexidade de Algoritmos;
\item Introdução a Árvores;
\item Árvores I;
\item Árvores de Busca II;
\item Hashing;
\item Gerência de Arquivos;
\item Ordenação;
\end{itemize}
\end{frame}

\begin{frame}{Bibliografia}
\begin{enumerate}
\item Material de Estruturas de Dados:
\begin{itemize}
\item Horowitz: Estruturas de Dados. Ed. Campus;
\item Wirth: Data Structures + Algorithms = Programs;
\item Outros livros, a maioria em inglês, todos disponíveis na BU ou no INE;
\end{itemize}
\item Material sobre  C++:
\begin{itemize}
\item Livros sobre Estruturas de Dados em “C” e C++;
\item Material sobre linguagem C++;
\end{itemize}
\item Bibliografia de Ferramentas de Programação:
\begin{itemize}
\item Livros sobre Unix / Linux;
\item Manuais de Editores;
\item Manuais de Programação;
\item Material Online.
\end{itemize}
\end{enumerate}
\end{frame}

\begin{frame}{Contato e Marcação de Horários}
\begin{itemize}
\item Monitoria: será disponibilizado um horário por semana para atendimento pelo estagiário de docência ou bolsista da disciplina;
\item Procedimento:
\begin{itemize}
\item Envie e-mail explicando seu problema;
\item Utilize sempre o Fórum de Discussão;
\item Se possível, o problema será resolvido por e-mail;
\item Caso contrário, será indicado o horário de atendimento.
\end{itemize}
\item Marcação de Horários
\begin{itemize}
\item Professores por mensagem no moodle;
\item Estagiário de docência e bolsista por e-mail divulgado no moodle;
\end{itemize}
\end{itemize}
\end{frame}

{
\usebackgroundtemplate{\includegraphics[width=\paperwidth,
height=\paperheight]{../reusable_images/fundo_capa.png}}
\begin{frame}

{\LARGE Perguntas????}

\end{frame}
}


{
\usebackgroundtemplate{\includegraphics[width=\paperwidth,
height=\paperheight]{../reusable_images/fundo_capa.png}}
\begin{frame}
\includegraphics[scale=0.8]{../reusable_images/cc_logo_arge.png}\hspace{0.5cm}
\includegraphics[scale=0.95]{../reusable_images/by.png}

\vspace{1cm}
Este trabalho está licenciado sob uma Licença Creative Commons Atribuição 4.0 Internacional. Para ver uma cópia desta licença, visite http://creativecommons.org/licenses/by/4.0/.

\end{frame}
}
\end{document}
